\documentclass[11pt]{article}
\usepackage[margin=2.5cm]{geometry}
\usepackage{amsmath,amssymb}
\usepackage{booktabs,array}
\usepackage{enumitem}
\usepackage[hidelinks]{hyperref}

\title{Response to ``Countdown: Faster Exact Search via Three-Way Joins''}
\author{A review prepared with Claude Opus~5.5 in Claude Code}
\date{10 October 2026}

\begin{document}
\maketitle

\begin{abstract}
The note under review, dated 10 October 2026 and written with ChatGPT, lowers the leading exponent for
deciding whether a Countdown target can be hit exactly from $\tfrac23$ to $0.666223$, in expected time
$2^{0.666223\, n \log_2 n + O(n)} \operatorname{poly}(B)$. We recomputed every numerical inequality it
relies on and checked the lemmas we could check without its main source. The arithmetic is right and the
reductions we examined are sound. The result is nevertheless much narrower than its title suggests. It
rests on a preprint published five days earlier that we could not inspect. It does not cover the problem
our solvers solve, which is the closest value with all four operations and a canonical answer. By the
note's own Proposition~5.1 it is beaten by the ordinary subset recursion whenever the card values are
bounded, as they are in the game. And for the six numbers of the game it changes the bound by a factor of
about $1.005$.
\end{abstract}

\section{What the note claims}

The note treats the number of input numbers $n$ as growing, with their total bit length $B$ a
polynomial factor. It splits every expression tree at a centroid node into two subtrees and the
\emph{context} above them, a path whose effect on the hole value is a fractional linear map
(Lemma~1.1). Inverting the context at the target turns the search into matching three lists, one entry
from each, under addition or multiplication. Lemma~2.1 turns multiplicative matching into seven
additive instances by rounding logarithms, so any three-list 3SUM algorithm in time $\tilde O(K^\tau)$
gives an exact solver in time $2^{\frac{\tau}{\tau+1} n \log_2 n + O(n)} \operatorname{poly}(B)$. With
$\tau = 2$ this is the earlier $\tfrac23$. Plugging in parameters from a 3SUM algorithm of Alman and
Vassilevska Williams~\cite{av} gives $\tau = 1.99657$ and the exponent $0.666285$. An asymmetric join
(Section~3) that uses the actual sizes of the three lists then gives $0.666223$.

\section{What we checked}

\paragraph{The constants.}
The repository script \texttt{paper/check\_fast\_solver.py} recomputes every numerical condition in
Sections~2 and~3. All of them hold.

\begin{center}
\small
\begin{tabular}{@{}llr@{}}
\toprule
Condition & Value & Margin \\ \midrule
$2\beta + \eta = \alpha$ & exact & --- \\
$\gamma > 0.06055$ & $0.0605550$ & $5.0 \times 10^{-6}$ \\
$q < 0.40916$ & $0.4091519$ & $8.1 \times 10^{-6}$ \\
$R > \alpha = 0.05668$ & $0.0566857$ & $5.7 \times 10^{-6}$ \\
$\alpha\gamma > \eta$ and $\beta - \alpha q > \eta$ & $0.0034320$, $0.0034338$ & $2$--$4 \times 10^{-6}$ \\
$\tau/(\tau+1) = 199657/299657$ & $0.6662851$ & exact \\
Lemma~3.1: $\gamma > 0.07735$, $q < 0.4549$ & $0.0773704$, $0.4548759$ & $2.0 \times 10^{-5}$ \\
Lemma~3.1: $\alpha\gamma > g$, $\alpha(1 - 2q) > h$ & $0.0039913$, $0.0046543$ & $1$--$4 \times 10^{-6}$ \\
Theorem~3.2: $(2 - g)/(3 - g) < K$ & $0.66622274$ & $5.7 \times 10^{-7}$ \\
Theorem~3.2: $(4 - 3\delta + h)/(6 - 4\delta + 2h) < K$ & $0.66621745$ & $5.6 \times 10^{-6}$ \\
\bottomrule
\end{tabular}
\end{center}

The margins are small but far above double-precision rounding, so the conditions are not in doubt. The
note says an accompanying script checks them with exact rational intervals; that script was not supplied.
One condition could not be checked: Lemma~3.1 also needs $R_{23}(\gamma) > \alpha$, and the note does not
give the formula for $R_{23}$.

\paragraph{Lemma 2.1.}
Correct. If $ab \ne c$ then $ab/c = P/R$ with distinct $P, R \le 2^{3b}$, so
$|\log(ab/c)| \ge 2^{-3b}$. With $Q = 2^{3b+4}$ this gap is at least 16, against a rounding error of at
most 3, so the seven shifted additive equalities capture exactly the multiplicative ones.

\paragraph{Theorem 3.2.}
The case analysis holds, but the case that needs the new asymmetric bound is a sliver. The largest of the
three parts, $a$, must satisfy $K/(2 - \delta) < a < 1 - K$, that is
$0.333684 < a < 0.333777$. Everywhere else the simpler pair or block joins already give $K$. The whole
improvement from $0.666285$ to $0.666223$ comes from that window.

\paragraph{Proposition 6.2.}
Sound. Encoding each vector as a product of distinct odd primes makes divisibility $y_b \mid x_a$
equivalent to orthogonality, and the factor $2^K$ rules out a legal division in the other direction.
Scaling both lists by $C$ puts every other one-operation result at or above $C$.

\section{Limitations}

\paragraph{The main step rests on a source we could not check.}
Both new exponents come from Lemmas~36 and~37 and Corollary~31 of~\cite{av}, an arXiv preprint dated
5 October 2026, five days before the note. We could not inspect it, so we cannot confirm that those
results say what the note uses. The note also states without proof that the preprint's algorithm for real
numbers can be simulated exactly on $b$-bit integers. Until both points are confirmed, $0.666223$ is a
conditional result.

\paragraph{It does not cover the problem the solvers solve.}
Our solvers return the value closest to the target, using all four operations, with a fixed canonical
answer. The note's fast bounds are for the exact-target question only. Its Section~6 says the
centroid method does not extend to closest values once exact division is allowed, and its Theorem~6.1
removes division to get a closest-value bound. For the full task the note's bound is its
Proposition~5.1, the subset recursion over all $2^n - 1$ supports, which is the algorithm the solvers
already use. Its canonical order also differs from ours: it orders the two children of a node by a
bit-mask of the cards they use, not by value.

\paragraph{For the real game the simple method wins.}
Proposition~5.1 shows that when every card is at most a fixed $M$, which is the case in Countdown
($M = 100$), the full closest-value problem takes $2^{O(n)} \operatorname{poly}(B)$ time. That grows more
slowly than $2^{0.666\, n \log_2 n}$. The new exponent only matters when card values can be arbitrarily
large.

\paragraph{The gain is negligible at six numbers.}
With $n = 6$, $n \log_2 n = 15.51$, so the change from $\tfrac23$ to $0.666223$ alters the leading term by
$2^{(2/3 - 0.666223) \cdot 15.51} \approx 1.005$. The 3SUM algorithms behind it rely on fast matrix
multiplication with very large constant factors, and the note says itself that it makes ``no claim of
practical speedup or research novelty''. Our fastest solvers already take about $0.012$\,ms per round.

\paragraph{No certificates for ``no solution''.}
The note says that no separate certificates for negative answers are considered. In our study every
unsolvable round has a machine-checked proof of that, through Lean and, for a pilot set, through SAT
proofs checked by a verified checker. A faster decision procedure without certificates gives weaker
evidence that a round cannot be solved.

\paragraph{Smaller points.}
The $\tfrac23$ bound it improves on comes from a ``previous note'' and a companion addendum~\cite{prev}
whose argument is quoted rather than shown. Lemma~3.1 and Theorem~3.2 would be easier to follow if the
asymmetric reduction of Chan, Vassilevska Williams and Xu~\cite{cvx} were restated, not only cited by
problem number.

\section{Conclusion}

As far as we could check, the note's mathematics is right: every numerical condition holds and the
reductions we examined are sound. Its result is an asymptotic upper bound for the exact-target problem with
unbounded card values. It depends on a preprint we could not inspect, and it does not apply to the
closest-value problem that Countdown solvers solve. For bounded cards the note itself shows that the
subset recursion is asymptotically better, and at six numbers the improvement is about half a percent of
the leading term. We would mention it in our paper only as a remark on the theory of the exact-target
problem, and would not quote $0.666223$ until~\cite{av} can be checked.

\begin{thebibliography}{9}
\bibitem{av} J.~Alman and V.~Vassilevska Williams. Truly subquadratic 3SUM and truly subcubic APSP via
triangles in sparse lopsided graphs. arXiv:2610.06783v1, 5 October 2026. Cited by the note; not inspected.
\bibitem{cvx} T.~M.~Chan, V.~Vassilevska Williams and Y.~Xu. Hardness for triangle problems under even
more believable hypotheses. arXiv:2203.08356v2, 2022.
\bibitem{prev} Countdown: hard bounds and a faster certifying solver. Companion addendum to the note,
10 October 2026.
\end{thebibliography}

\end{document}
